% TOP_PROBLEM: {"id": 2306027, "title": "Research Problems in Function Theory — Problem 6.27", "category_id": 8, "category": {"id": 8, "name": "analysis", "display_name": "Analysis", "description": "Limits, continuity, calculus, and function theory.", "slug": "analysis", "order_index": 8, "created_at": "2026-07-31T15:26:25.671Z"}, "status": "open"}
% FIRST_POSTED: null
% ENTRY_KIND: statement_only
% Imported from: batch02.tex; UnsolvedMath ID 2306027
% Compile twice: lualatex 2306027.tex
\documentclass[11pt]{article}
\usepackage{fontspec}
\setmainfont{Latin Modern Roman}
\usepackage{amsmath,amssymb,mathtools,unicode-math}
\setmathfont{Latin Modern Math}
\usepackage{xurl,hyperref}
\hypersetup{hidelinks,pdftitle={UnsolvedMath Problem 2306027}}
\newcommand{\sref}[2]{\hyperlink{source:#1:#2}{[S#2]}}
\newcommand{\cataloguescope}[1]{\paragraph{Mathematical question.}#1}
\begin{document}
% UnsolvedMath ID 2306027; source catalogue code AMR-022-6027
\setcounter{section}{2306026}
\section{Research Problems in Function Theory — Problem 6.27}\label{2306027}
{\small\sffamily UnsolvedMath ID 2306027 · Analysis}
\subsection{Definitions and mathematical statement}

\paragraph{Shared notation.}$\mathbb N=\{1,2,\ldots\}$, and $\mathbb Z,\mathbb Q,\mathbb R,\mathbb C$ denote the integers, rationals, reals and complex numbers. Any other conventions are specified locally.


Let $g$ be holomorphic and injective on $\{|z|>1\}$, with Laurent expansion $g(z)=z+b_0+\sum_{n\ge1}b_nz^{-n}$. For integers $n\ge3$, define $H_n=\max\{\nu(|b_\nu|+1):\nu\ge1,\ 0<|\nu-n|<n/2\}$. This is the literal coefficient expression recorded in the source. The exterior area theorem gives $|b_n|\le n^{-1/2}$; also $H_n\ge n/2$, so this literal version already has an elementary bound.
\paragraph{Statement recorded in the supplied source.}
For every $\varepsilon>0$, does there exist $C_{g,\varepsilon}<\infty$ such that $n|b_n|\le C_{g,\varepsilon}n^\varepsilon H_n$ for all sufficiently large $n$? \sref{2306027}{2}
\subsection{Short English statement}
Record the exterior coefficient comparison as printed, including its neighboring-coefficient expression. That literal version follows from the area theorem.
\subsection{Sources}
\begin{description}
\item[\hypertarget{source:2306027:1}{S1}] ULAM.AI and the UnsolvedMath Contributors, UnsolvedMath: A Curated Collection of Open Mathematics Problems, record 2306027 (AMR-022-6027). CC BY 4.0. \url{https://huggingface.co/datasets/ulamai/UnsolvedMath}.
\item[\hypertarget{source:2306027:2}{S2}] W. K. Hayman and E. F. Lingham, Research Problems in Function Theory (2018), arXiv:1809.07200, Chapter 6, Problem 6.27, complete update and preceding notation. \url{https://arxiv.org/abs/1809.07200}.
\end{description}
\label{end:2306027}
\end{document}
